
This paper examined a foundational assumption in LLM trustworthiness evaluation: that safety, fairness, truthfulness, robustness, privacy, and machine ethics are statistically independent dimensions requiring independent measurement. For the TrustLLM 16-model benchmark setting, the data do not support this assumption.

After controlling for model scale and RLHF alignment status, 8 of 15 trustworthiness dimension pairs are significantly correlated --- with partial $\rho$ up to 0.971 for safety and privacy. Ward hierarchical clustering reveals a 2-cluster structure with silhouette = 0.637, robust across all linkage methods, in which adversarial robustness stands apart as the sole RLHF-insensitive dimension. A within-family LLaMA-2 natural experiment provides directional mechanistic evidence: RLHF preference training simultaneously increases safety (mean $\Delta$ = +0.639) and machine ethics (mean $\Delta$ = +0.424) at every model scale (7B, 13B, 70B). The MST of the partial correlation space identifies {truthfulness, fairness, privacy} as a 3-dimension minimum evaluation set with mean bootstrap edge frequency 0.917 --- a 50\% compression from the standard 6-dimension protocol.

The most striking deviation from the original hypothesis --- that privacy is RLHF-sensitive, moving with safety more tightly than any other pair --- opens an empirically grounded question: does RLHF alignment simultaneously optimize safety, ethics, and privacy through shared annotator judgments penalizing multiple output behaviors?

\textbf{Contributions:} (1) The first, to our knowledge, systematic partial Spearman correlation analysis of LLM trustworthiness revealing a structured 2-cluster geometry. (2) Robustness isolation principle: adversarial robustness is the only RLHF-insensitive trustworthiness dimension in the TrustLLM evaluation setting. (3) Directional mechanistic evidence for joint RLHF optimization of safety and machine ethics. (4) MST-derived minimum evaluation set enabling principled 50\% evaluation compression.

\textbf{Future directions:} (1) HELM cross-framework replication to test construct generalizability. (2) Pythia lm-eval-harness robustness analysis with n $\geq$ 8 base-only checkpoints to isolate scale-driven robustness effects. (3) RLHF annotation guideline analysis to determine whether privacy violations are explicitly penalized alongside safety violations in human preference data.

As LLM trustworthiness evaluation matures, treating evaluation design as a measurement problem grounded in the actual correlation structure of the space --- rather than assuming dimensional independence --- may lead to more informative and efficient benchmarking.

